\chapter{Introducción a Bloques y Capas}

\section{Gestión Eficiente con Bloques y Capas}

La eficiencia en el dibujo CAD no solo depende de la velocidad de
dibujo, sino de cómo se gestiona y organiza la información. El uso
adecuado de \textbf{bloques} y \textbf{capas} es lo que diferencia a un
dibujante aficionado de un profesional.

\subsection{El Poder de los Bloques}

Un bloque es un conjunto de entidades geométricas (líneas, arcos,
círculos, texto) que se agrupan para formar una sola entidad compleja.

\textbf{Ventajas principales:}
\begin{enumerate}
	\item \textbf{Reutilización:} Permite insertar elementos estándar
	      (ej: tornillos, puertas, símbolos eléctricos) repetidas veces
	      en un mismo dibujo o en dibujos diferentes, ahorrando tiempo
	      y espacio.
	\item \textbf{Actualización global:} Si el bloque se edita, todas
	      las instancias insertadas en el dibujo se actualizan
	      automáticamente.
	\item \textbf{Reducción del tamaño de archivo:} Al insertar múltiples
	      copias de un bloque, el archivo ocupa mucho menos que si
	      se dibujara cada elemento individualmente.
\end{enumerate}

\subsection{La Organización mediante Capas (Layers)}

Las capas funcionan como láminas de transparencia superpuestas. Cada
capa permite gestionar la información del dibujo de forma lógica y
ordenada.

Cada capa tiene propiedades que definen el comportamiento de todas las
entidades dibujadas en ella:
\begin{itemize}
	\item \textbf{Visibilidad (On/Off):} Permite ocultar elementos
	      temporales.
	\item \textbf{Congelado (Freeze):} Similar a apagar, pero además
	      la computadora no procesa los datos, aumentando la velocidad.
	\item \textbf{Bloqueo (Lock):} Impide editar elementos de esa capa
	      sin ocultarlos.
	\item \textbf{Propiedades:} Color, tipo de línea (continua,
	      discontinua, eje) y grosor de línea (fundamental para la
	      calidad en la impresión).
\end{itemize}

La correcta gestión de capas es vital para imprimir planos técnicos,
donde se necesita controlar qué información es visible y cuál debe
tener un grosor de línea específico (ej: trazos gruesos para contornos y
finos para cotas).

\section{Cuestionario}
\begin{enumerate}
	\item Define bloque y explica su principal ventaja en términos de
	      actualización.
	\item ¿Cómo se edita un bloque después de haber sido insertado?
	\item ¿Para qué sirve el gestor de capas en proyectos grandes?
	\item Diferencia técnica entre un bloque y un grupo de entidades.
	\item ¿Cómo ayuda el uso de capas en el proceso de impresión?
	\item ¿Qué sucede al "congelar" una capa y por qué mejora el
	      rendimiento?
	\item ¿Qué es un atributo de bloque y para qué sirve?
	\item ¿Para qué se utilizan los colores asignados a las capas?
\end{enumerate}

\section{Parte Práctica}
\begin{enumerate}
	\item Crear un bloque que represente una puerta normalizada y
	      guardarlo en la biblioteca.
	\item Insertar varias copias del bloque en un plano de planta complejo.
	\item Crear un bloque con atributos para el número de habitación (ID).
	\item Organizar un plano en al menos 5 capas (ej: Muros, Muebles,
	      Cotas, Texto, Ejes) con colores y tipos de línea distintos.
	\item Cambiar el color de todas las entidades en la capa "Muebles"
	      rápidamente.
	\item Bloquear la capa "Muros" para no editar accidentalmente las
	      paredes mientras se dibuja la instalación eléctrica.
	\item Crear un bloque anidado (un bloque dentro de otro).
	\item Cambiar el tipo de línea de la capa "Ejes" a "línea trazo-punto".
	\item Insertar un bloque existente desde otro archivo de dibujo.
	\item Gestionar la visibilidad de las capas para imprimir solo el
	      plano de arquitectura, ocultando instalaciones.
\end{enumerate}
